\section{The hypotheses that were explored and rejected}
\label{sec:rejected}

Five candidate architectures were built in isolated worktrees, gated before they were scored, and
certified against the side-channel definition by the harness of \S\ref{sec:methodology-side} rather
than by inspection. \textbf{One survived, and it is the configuration change of
\S\ref{sec:configuration} rather than an architecture.} The four that died are set out here with the
measurements that killed them, because each closes a standing entry in the project's design register
with a number instead of an argument, and because a closed direction is the part of this cycle that
does not depreciate.

Every figure in this section is on \emph{training} material, at the budget stated with it, and none
of it is certified: every pooled interval reported below has a lower bound under the
\vsevenPrimaryFloor{}-point floor of \S\ref{sec:methodology-floor}. They are measured edges, and the
distinction is kept.

\subsection{The declaration-allocation channel closes at exactly zero}
\label{sec:rejected-alloc}

The register's second-ranked entry, priority \vsevenLonePriority{}, held that the whole of \vsix{}'s
margin lives in the declaration channel and that roughly two points of it remained. Its named
mechanism was that the deployed allocator --- a product of marginals --- picks a different allocation
from the joint posterior's argmax often enough to matter.

\textbf{It does not.} Over $\sim$\vsevenLoneDecls{} recorded voluntary declarations, across two banks
and both urgency arms, the joint argmax and the marginal-product argmax select \emph{the same
allocation at every one}, with the joint score holding a strict opinion on \vsevenLoneStrictLo\%--\vsevenLoneStrictHi\% of genuinely ambiguous cases. The reason is algebraic rather than statistical: the sequential joint's first
chain-rule factor \emph{is} the marginal, and the re-normalisation that follows reweights the
survivors without reordering them.

Two further measurements close the direction rather than merely the mechanism. First, the
\emph{exact} joint maximiser is a worse decision rule than the deployed approximation on the
declaration decision itself: paired, it fixes \vsevenJointFixed{} declarations and breaks \vsevenJointBroken{}, a net
\vsevenJointWorseNet{}~pp, which converts to about \vsevenJointWorseWin{}~win-rate points. This is the
second independent instance in this corpus of the project's standing habit paying off --- ``exact
Bayesian inference'' is an assumption to test, not a property to assume --- the first being \vsix{}'s
finding that the exact posterior is a worse predictor of card location than the approximation the
policy deploys. Second, \vsevenJointFlatLo\%--\vsevenJointFlatHi\% of the error class sits in
flat-posterior states where the exact object is provably a coin flip, so most of what looked like
headroom is an information limit and not a policy defect.

\textbf{And the channel is largely downstream of a defect that the configuration change removes
anyway.} Deleting the urgency escalation removes \vsevenLoneUrgencyShare\% of the entire
allocation-error class and drops the oracle ceiling for a \emph{perfect} allocator from
\vsevenLoneCeilingBefore{} win-rate points to \vsevenLoneCeilingAfter{} --- below the detection floor.
The consequence is a composition constraint that this paper honours: a declaration-channel gain and
the urgency-off gain may never be added, because they are the same defect counted twice.

\subsection{A leaf evaluator fitted \vsevenKoneFold$\times$ better buys nothing}
\label{sec:rejected-leaf}

The inherited conditional held that the test-time search cannot be judged until its leaf evaluator is
rebuilt, because the shipped one is algebraically close to a rescaling of the hit probability. A leaf
was fitted for the purpose, and the statistic that matters --- \emph{between-candidate} $R^2$, the
only component a paired deviation rule can consume --- improved from \vsevenKoneRsqBefore{} to
\vsevenKoneRsqAfter{}, a \vsevenKoneFold$\times$ gain replicated out of sample on $\sim$\vsevenKoneLeaves{} leaves.

In play it is worth \vsevenKoneInPlay{} against the material control's \vsevenKoneControl{}, losing on
both banks individually. The diagnosis is conversion failure, not fit failure: the rule's deviation
rate is \vsevenKoneChangeLo{} for the material leaf and \vsevenKoneChangeHi{} for the retuned one, so
the deviation rule at the shipped shrinkage is \textbf{nearly leaf-invariant}. A better leaf is not
the lever.

The conditional itself is discharged separately, and by re-measurement rather than by the new
evaluator. The single underpowered cell the conditional rested on --- one \vsevenFullGameOldN{}-game cell at \vsevenFullGameOld{} --- re-runs at \vsevenTruncatedEndgame{} replicated in the endgame regime with
\vsix{}'s \emph{own} leaf. Full-game truncated search already works with the leaf the conclusion said
could not support it; it is simply at a dominated operating point.

\subsection{Fitting against a per-decision objective produces a worse policy}
\label{sec:rejected-perdecision}

The register's enabling entry held that a per-decision objective is necessary because the scoreboard
is too noisy to resolve small effects. \textbf{Its precision claim is confirmed and its fitting form
is killed}, and the two halves are independent.

The design effect of deal clustering was measured for the first time in this cycle, and a
one-point-equivalent effect resolves in \vsevenPerDecisionGamesLo{}--\vsevenPerDecisionGamesHi{} games
on the declaration channel against \vsevenScoreboardGames{} on the scoreboard: the precision claim
confirmed at \vsevenDesignEffectLo$\times$--\vsevenDesignEffectHi$\times$.

But all three self-oriented per-decision objectives, fitted at matched budget against a per-game
control, produced \emph{worse} policies: \vsevenKfourSelfDecl{}, \vsevenKfourSelfAsk{} and
\vsevenKfourSelfAlloc{}~points, every one replicated in sign and two of them clearing the floor as
certified negatives. Each fit moved its own proxy in the intended direction and lost anyway:
the declaration-oriented fit bought \vsevenKfourTradeGain{}~pp of declaration accuracy by paying
\vsevenKfourTradeCost{}~pp of ask accuracy, and the ask-oriented fit did the reverse. The two
per-decision channels trade against each other, so a fitter holding \vsevenCoordinates{} coordinates buys one out of the other. Keep the estimator; drop the objective.

\subsection{Nothing learned survives a free random tie-break}
\label{sec:rejected-tie}

\vsevenTieContested\% of contested ask decisions in \vsix{} tie bit-for-bit. The open question was
whether anything can be \emph{selected} inside that group, as opposed to merely randomised. A learned
re-ranker was built --- the corpus's first learned agent --- restricted to the tie group.

Against the deployed policy it scores \vsevenKfiveTieOnly{} pooled, which is indistinguishable from the free hash tie-break's \vsevenKfiveFreeTie{}; head to head against that free tie-break at \vsevenPhaseTwoCell{} games it is
\vsevenKfiveHeadToHead{} $\pm$ \vsevenKfivePm{}, replicated. \textbf{Nothing learned survives the free
random tie-break.} The correct control for anything operating inside the tie group is the
decorrelated policy and not the deployed one, which is a methodological point this paper inherits
and applies.

The same capture explains why. Setting every candidate's true advantage to exactly zero and feeding
the capture's real standard errors through the engine's real deviation rule reproduces the search's observed deviation rate to a \vsevenCurseRelErr\% relative error on a quantity the null was never fitted to: the
search's per-decision selection signal is $\sim$\vsevenCurseShare\% winner's curse on its own
Monte-Carlo draw, and $\sim$\vsevenCurseTieShare\% inside the tie group, where the rule applies no
shrinkage at all.

\subsection{Two directions closed on arithmetic before they were built}
\label{sec:rejected-arithmetic}

\paragraph{Forced-endgame allocation.} The register priced closing the entire forced-endgame accuracy
gap. Re-derived from this cycle's own capture it is worth \vsevenLthirteenClose{}~win-rate points ---
the register's own figure corrected \emph{upward} by a factor of two and still a fiftieth of the
detection floor. The one route by which it could have returned, an adversary driving incidence up, was measured and costs the adversary \vsevenLthirteenAdvCost{}~points.

\paragraph{A stochastic policy bought to reduce readability.} The brief made this conditional on the
white-box inverter biting. Transcript inversion contracts the deal posterior by $\sim$\vsevenBitsPerAsk{} bits per observed
ask --- the first such measurement in this corpus, and one the threat model's literature review
found no external precedent for --- and
converting those bits into play is worth \vsevenInversionUnfitted{}~points unfitted. But the excess
tracks planted-edge size within $\pm$\vsevenInversionTrack{} on all eight rungs, which is the
signature of a general strength gain and not of an exploitation instrument. And no linear handicap
adds more than \vsevenHandicapBits{}~bits per ask, so within this policy family \textbf{there is no
readability handicap for a stochastic policy to buy.} The antecedent fails and the limb is closed
rather than skipped.

\subsection{What the four kills leave}
\label{sec:rejected-leaves}

The register the cycle inherited described a policy with several points of accessible headroom in
four named places. Three of them are now closed with measurements and the fourth is confirmed as an
instrument rather than a lever. What remains as the located gain is a single coordinate found by an
unfitted sweep --- not by any of the architecture work --- and \S\ref{sec:results-attribution}
measures it on holdout. That is the barrier this paper names, and \S\ref{sec:conclusion} states it
as such.
